\documentclass[11pt]{article}
\usepackage[a4paper,margin=1in]{geometry}
\usepackage{amsmath,amssymb,mathtools,booktabs}
\title{Finite harmonic expansion of the evaluated-phase residual flux}
\author{}
\date{2026-09-09}
\begin{document}
\maketitle
\small
\section{Scope and chart}
This memo expands the evaluated-phase coefficient already defined in
\texttt{moments\_body.tex} and the finite chart of
\texttt{quantitative\_extension.tex}.  It is a fixed-band, finite-label,
compact-support statement.  The physical viscosity is the positive number
$\nu_{\rm NS}$; it is not replaced by $\varepsilon$ or by the source clock
$\nu=\sqrt{A_0}$.  Put
\[
 0<Q\leq1,\qquad A=\tfrac12+h,\qquad D=\tfrac12-h,
 \qquad \varepsilon=Q^h,
 \qquad r=\sqrt Q R,\quad z=Q^D Z.
\]
The source constants $A_0,\lambda_0$ and $\mu=\lambda_0$ retain their given
values.  Every statement below is at one such $Q$ and on the strict interior
compact patch where all retained coefficients and their displayed
$Z$-derivatives are smooth and compactly supported in $R$.

\section{Exact angular kernel}
Let $I_*$ and $I_y$ be the finite retained old-parent and new-harmonic label
sets.  For $\tau\in I_*$ and $j\in I_y$, let $n_\tau,n_j\in\mathbb Z\setminus\{0\}$
be their actual angular carriers.  Write their cylindrical component
coefficients, before taking real parts, as
\[
 W_{\tau,\ell}=\Re(a_{\tau,\ell}e^{in_\tau\theta}),\qquad
 V_{j,\ell}=Q^{-A}\Re(\zeta_{j,\ell}e^{in_j\theta}),
 \qquad \zeta_{j,\ell}=b_{j,\ell}e^{i\varphi_j},
\]
where $\ell\in\{r,\theta,z\}$ and
$\varphi_j=k_{\beta_j}m_j\Phi_j$.  The old coefficient $a_{\tau,\ell}$ includes its full retained old phase and all cutoff/path factors, so its $Z$ derivative is the full old axial derivative in this representation.  Thus
\[
 \partial_Z\zeta_{j,\ell}=\delta_{j,\ell}e^{i\varphi_j},
 \qquad
 \delta_{j,\ell}:=\partial_Zb_{j,\ell}
       +ik_{\beta_j}m_j(\partial_Z\Phi_j)b_{j,\ell}.
\]
For complex $x,y$ and integers $n,m\ne0$ define
\[
 \mathcal A_{n,m}(x,y):=\frac14\left[
 {\bf1}_{n+m=0}(xy+\bar x\bar y)
 +{\bf1}_{n-m=0}(x\bar y+\bar x y)\right].
 \tag{E1}\label{eflux:kernel}
\]
Expanding the two real parts and integrating each exponential over
$[0,2\pi]$ proves exactly
\[
 \left\langle\Re(xe^{in\theta})\Re(ye^{im\theta})\right\rangle_\theta
 =\mathcal A_{n,m}(x,y).
 \tag{E2}\label{eflux:angular}
\]
No equal-frequency assumption is made: the indicators in
\eqref{eflux:kernel} retain both $n=m$ and $n=-m$ possibilities.  Moreover,
for real $Z$, differentiation commutes with the kernel,
\[
 \partial_Z\mathcal A_{n,m}(x,y)=
 \mathcal A_{n,m}(\partial_Zx,y)+\mathcal A_{n,m}(x,\partial_Zy),
 \tag{E3}\label{eflux:kernel-derivative}
\]
because complex conjugation and the integer indicators commute with
$\partial_Z$.

\section{Finite stress expansion}
Let $\Pi_Y f=\int_{\mathbb T^2}f\,dY$ and $\mathsf R_Y=I-\Pi_Y$.
The stress in the moments calculation is
$\mathcal S=\langle w_*\otimes v+v\otimes w_*+v\otimes v\rangle_\theta$,
and $\mathcal S^\circ=\mathsf R_Y\mathcal S$.  Since
$S^*=Q^{2A}S$, the exact $z,\tau_e$ component is
\begin{align}
 (S^\circ_{z,\tau_e})^*
 =\mathsf R_Y\Bigg\{&Q^A\sum_{\tau\in I_*}\sum_{j\in I_y}
 \left[\mathcal A_{n_\tau,n_j}(a_{\tau,z},\zeta_{j,\tau_e})
       +\mathcal A_{n_j,n_\tau}(\zeta_{j,z},a_{\tau,\tau_e})\right]
 \nonumber\\[-2pt]
 &+\sum_{i,j\in I_y}
       \mathcal A_{n_i,n_j}(\zeta_{i,z},\zeta_{j,\tau_e})\Bigg\}.
 \tag{E4}\label{eflux:stress}
\end{align}
The first line is every old--new and new--old product; the second is the
complete new quadratic sum, including both ordered pairs.  Terms whose
angular carriers do not satisfy an indicator in \eqref{eflux:kernel} vanish
under the angular average, while their labels remain available for the
finite sum.  The $Y$-projection is applied after the angular product, so
$\langle\mathcal S^\circ\rangle_Y=0$ is retained without cancelling the
physical evaluation.

Applying \eqref{eflux:kernel-derivative} gives the exact differentiated
coefficient.  Define
\[
 \mathcal D_{n,m}(x,y):=\mathcal A_{n,m}(\partial_Zx,y)
                       +\mathcal A_{n,m}(x,\partial_Zy).
\]
Then
\begin{align}
 \partial_Z(S^\circ_{z,\tau_e})^*=\mathsf R_Y\Bigg\{&Q^A\sum_{\tau,j}
 \big[\mathcal D_{n_\tau,n_j}(a_{\tau,z},\zeta_{j,\tau_e})
      +\mathcal D_{n_j,n_\tau}(\zeta_{j,z},a_{\tau,\tau_e})\big]
 \nonumber\\
 &+\sum_{i,j}\mathcal D_{n_i,n_j}(\zeta_{i,z},\zeta_{j,\tau_e})\Bigg\},
 \tag{E5}\label{eflux:stress-derivative}
\end{align}
where all sums have the finite ranges $\tau\in I_*$ and $i,j\in I_y$.
In \eqref{eflux:stress-derivative}, the old terms contain
$\partial_Za_{\tau,\ell}$ and the new terms contain
$\partial_Z\zeta_{j,\ell}=\delta_{j,\ell}e^{i\varphi_j}$; no parent derivative
has been identified with a new derivative.

\section{The coefficient $C_e^\circ$ and its finite envelope}
For $e\in\{1,2\}$, set $\tau_2=\theta$ and $\tau_1=z$, and define exactly
\[
 C_e^\circ(Z,T):=\int_0^\infty R^e\sup_{Y\in\mathbb T^2}
 \left|\partial_Z(S^\circ_{z,\tau_e})^*(R,Z,T,Y)\right|\,dR.
 \tag{E6}\label{eflux:C}
\]
Because $\|\mathsf R_Yf\|_{L^\infty_Y}\le2\|f\|_{L^\infty_Y}$, the finite sums in \eqref{eflux:stress-derivative} give the explicit envelope
\begin{align}
 C_e^\circ\le 2Q^A\sum_{\tau,j}\int_0^\infty R^e\sup_Y &\left(
 |\mathcal D_{n_\tau,n_j}(a_{\tau,z},\zeta_{j,\tau_e})|
 +|\mathcal D_{n_j,n_\tau}(\zeta_{j,z},a_{\tau,\tau_e})|\right)dR
 \nonumber\\
 &+2\sum_{i,j}\int_0^\infty R^e\sup_Y
 |\mathcal D_{n_i,n_j}(\zeta_{i,z},\zeta_{j,\tau_e})|\,dR .
 \tag{E6a}\label{eflux:C-envelope}
\end{align}
The indicators in $\mathcal D$ make every nonmatching pair exactly zero;
therefore \eqref{eflux:C-envelope} is a finite pair sum, not an infinite
or continuum estimate.\nThe $Y$-evaluation map $Y=Y(r,t)$ is independent of $z$, hence the chain
rule in the moments body gives
\[
 M_e^{\rm ev}-M_e
 =\partial_z\int_0^\infty r^e\mathcal S^\circ_{z,\tau_e}(r,z,t,Y(r,t))\,dr,
 \qquad
 |M_e^{\rm ev}-M_e|\le
 Q^{(e+1)/2-2A-D}C_e^\circ
 =Q^{e/2-2A+h}C_e^\circ.
 \tag{E7}\label{eflux:evaluation-bound}
\]
The exponent uses $1/2-D=h$ exactly.

For a directly checkable finite majorant, let
\[
 \|f\|_{e,\infty Y}:=\left(\int_0^\infty R^e\sup_Y|f(R,Y)|^2dR\right)^{1/2},
 \quad \alpha_e:=\frac{e+1}{4}-A,
\]
and use the finite chart quantities
\[
 \|a_{\tau,\ell}\|_{e,\infty Y}\le Q^{\omega_{e,\tau}-(e+1)/4}W_{e,\tau,\ell},\quad
 \|\partial_za_{\tau,\ell}\|_{e,\infty Y}\le Q^{\dot\omega_{e,\tau}-(e+1)/4}\dot W_{e,\tau,\ell},
 \quad
 \|b_{j,\ell}\|_e\le B_{j,e,\ell},\quad
 \|\delta_{j,\ell}\|_e\le D_{j,e,\ell}.
 \tag{E8}\label{eflux:orders}
\]
These are finite compact-patch constants; the displayed chart powers include the physical-to-chart measure factor and do not assert uniformity in $Q$.
Since $\partial_Za=Q^D\partial_za$, Cauchy--Schwarz and $\|\mathsf R_Yf\|_\infty\le2\|f\|_\infty$
give the old--new contributions to the right side of \eqref{eflux:evaluation-bound}
with powers
\[
 Q^{\alpha_e+\dot\omega_{e,\tau}}\dot W_{e,\tau,\ell}B_{j,e,k},
 \qquad
 Q^{\alpha_e+\omega_{e,\tau}-D}W_{e,\tau,\ell}D_{j,e,k},
 \tag{E9}\label{eflux:old-orders}
\]
for every angularly matched pair and the appropriate components $(\ell,k)$.
The two exponents are distinct and are not merged unless the actual parent
proves $\dot\omega_{e,\tau}=\omega_{e,\tau}-D$.  The quadratic evaluated terms
carry the separate exact power
\[
 Q^{(e+1)/2-2A-D}=Q^{e/2-2A+h}
\tag{E10}\label{eflux:quad-order}
\]
against the finite products $D_{i,e,z}B_{j,e,\tau_e}$ and
$B_{i,e,z}D_{j,e,\tau_e}$.  Thus evaluation preserves the same finite
$Q$ powers as the moment rows after the explicit $C_e^\circ$ factor, while
retaining all harmonic matching and all old-parent derivative orders.

\section{Support, signs, and viscosity}
Every term in \eqref{eflux:stress} contains at least one compactly supported
new factor, so the $R$ integrals in \eqref{eflux:C} and the boundary terms in
the moments are finite and the radial boundary contribution vanishes.  The
kernel is real and retains the signs of the signed amplitudes and phases.
Full reflection reverses the corresponding signed cross flux as dictated by
the reflected vector field; absolute envelopes use the same finite constants.
No sign of $Q$, $\varepsilon$, $A$, $D$, or $\nu_{\rm NS}$ is changed.
The retained viscous term remains $-\nu_{\rm NS}\Delta v$ with its radial,
axial, angular, and cylindrical coupling costs.  Nothing in the evaluation
map supplies a source endpoint or a uniform $Q\downarrow0$ estimate.

\section{Finite conclusion}
Equations \eqref{eflux:stress}--\eqref{eflux:quad-order} are the strongest
finite evaluated-phase expansion established here: an exact finite harmonic
identity, its exact $Z$ derivative, and a finite compact-patch envelope with
separate old field and old derivative orders.  They do not prove the released
third shooting return, a uniform infinite modified sequence, or a smooth
terminal-force extension.
\end{document}







